% !TeX program = pdflatex
\documentclass[12pt,a4paper]{article}
\usepackage[T2A]{fontenc}
\usepackage[utf8]{inputenc}
\usepackage[russian]{babel}
\usepackage{amsmath,amssymb}
\usepackage[left=21mm,right=21mm,top=14mm,bottom=14mm,includeheadfoot,
  headheight=15pt,headsep=6mm,footskip=9mm]{geometry}
\usepackage{needspace,fancyhdr,tikz,tcolorbox}
\tcbuselibrary{skins,breakable}
\usepackage[unicode,hidelinks]{hyperref}
\hypersetup{pdftitle={Зачёт по КР2},pdfsubject={Дифференциальные уравнения. Ответы на 17 теоретических вопросов}}

\definecolor{navy}{HTML}{204F79}
\definecolor{leaf}{HTML}{2D795A}
\definecolor{slate}{HTML}{496DA4}
\definecolor{muted}{HTML}{626E7D}
\definecolor{gold}{HTML}{E3A735}
\definecolor{ochre}{HTML}{B86619}

\setlength{\parindent}{0pt}
\setlength{\parskip}{4pt plus 1pt}
\setlength{\emergencystretch}{2em}
\raggedbottom
\pagestyle{fancy}
\fancyhf{}
\fancyhead[L]{\small\color{muted}Дифференциальные уравнения}
\fancyhead[R]{\small\color{muted}Зачёт по КР2}
\fancyfoot[C]{\small\bfseries\color{navy}\thepage}
\renewcommand{\headrulewidth}{0pt}
\fancypagestyle{opening}{\fancyhf{}\fancyfoot[C]{\small\bfseries\color{navy}\thepage}}

\newcommand{\lectureref}[2]{%
  \href{https://github.com/alwaysnervous/diffury/blob/c23cbfa81e0849fda70a8fda2097cfeefa5db213/main/main.tex\#L#2}{\textcolor{navy}{\S\,#1}}}
\newcommand{\question}[4]{%
  \par\addvspace{12pt}\Needspace{8\baselineskip}%
  \pdfbookmark[1]{#1. #2}{question.#1}%
  \noindent\begin{minipage}{\linewidth}
    \tikz[baseline=(n.base)]\node[circle,fill=navy,text=white,inner sep=0pt,
      minimum size=7.2mm,font=\bfseries] (n) {#1};\hspace{2mm}%
    {\large\bfseries\color{navy}#2}\par
    \vspace{3pt}{\color{navy!28}\hrule height .6pt}\vspace{4pt}
    {\small\itshape\raggedright #3\par}
    {\small\color{muted}Основной конспект, #4\par}
  \end{minipage}\par\nobreak}

\tcbset{card/.style={enhanced,breakable,arc=2pt,boxrule=.5pt,
  boxsep=0pt,left=8pt,right=8pt,top=7pt,bottom=5pt,before skip=8pt,after skip=4pt,
  fonttitle=\small\bfseries,coltitle=white,
  attach boxed title to top left={xshift=3mm,yshift=-2mm},
  boxed title style={frame hidden,arc=1.5pt,boxrule=0pt,left=5pt,right=5pt,top=1pt,bottom=1pt},
  before upper={\setlength{\parskip}{3pt}\setlength{\abovedisplayskip}{5pt}%
    \setlength{\belowdisplayskip}{4pt}\setlength{\abovedisplayshortskip}{3pt}%
    \setlength{\belowdisplayshortskip}{3pt}}}}
\newcommand{\declarecard}[2]{%
  \newtcolorbox{#1}[1]{card,colback=#2!4!white,colframe=#2!80,
    colbacktitle=#2,title={##1}}}
\declarecard{definition}{leaf}
\declarecard{result}{navy}
\declarecard{lemma}{slate}
\declarecard{method}{navy}
\declarecard{note}{muted}
\declarecard{example}{ochre}

\newcommand{\idea}[1]{%
  \begin{tcolorbox}[enhanced,frame hidden,interior hidden,sharp corners,
    borderline west={1.5pt}{0pt}{muted!45},boxsep=0pt,left=7pt,right=0pt,
    top=1pt,bottom=1pt,before skip=5pt,after skip=5pt]
    {\small\bfseries\color{muted}Идея доказательства.}\enspace #1
  \end{tcolorbox}}

\begin{document}
\thispagestyle{opening}
\begin{tcolorbox}[enhanced,sharp corners,frame hidden,colback=navy,colupper=white,
  borderline south={2pt}{0pt}{gold},left=11pt,right=11pt,top=7pt,bottom=8pt,
  before skip=0pt,after skip=9pt]
  {\small Дифференциальные уравнения\par}
  {\bfseries\LARGE Зачёт по КР2\par}
  Ответы на 17 теоретических вопросов
\end{tcolorbox}
\question{1}{Уравнения с разделяющимися переменными}{Сформулируйте и поясните теорему существования и единственности
для уравнений с разделяющимися переменными. Расскажите о
частных случаях}
{\lectureref{9}{1430-L1462}, \lectureref{9.1}{1465-L1476}, \lectureref{9.2}{1566-L1577}}
\begin{result}{Теорема}
Пусть \(f\) непрерывна на интервале \(I\), а \(g\) непрерывна и отлична от нуля на интервале \(K\). Для каждой точки \((x_0,y_0)\in I\times K\) задача \(y'=f(x)g(y)\), \(y(x_0)=y_0\), имеет единственное локальное решение. Его задаёт равенство
\[
 \int_{y_0}^{y(x)}\frac{du}{g(u)}=\int_{x_0}^{x} f(t)\,dt .
\]
\end{result}
\idea{Делим на \(g(y)\) и интегрируем от начальной точки, получая \(G(y)-G(y_0)=F(x)-F(x_0)\), где \(G'=1/g\) и \(F'=f\). Поскольку \(1/g\) сохраняет знак, \(G\) строго монотонна, и формула \(y(x)=G^{-1}\!\bigl(F(x)+G(y_0)-F(x_0)\bigr)\) даёт единственное локальное решение.}

\textbf{При \(y'=f(x)\)} решение равно \(y_0+\int_{x_0}^x f(t)\,dt\), и кривые отличаются вертикальным сдвигом. Пусть \(f(x)\to\infty\) у конечной границы \(a_1\). Если интеграл до \(a_1\) сходится, кривая подходит к конечной точке прямой \(x=a_1\), где в расширенном поле направлений единственность может нарушиться. Если интеграл расходится, \(x=a_1\) --- вертикальная асимптота.

\textbf{При \(y'=g(y)\)} на интервале, где \(g\ne0\), получаем \(x=x_0+\int_{y_0}^y du/g(u)\). Кривые отличаются горизонтальным сдвигом. Если \(g(c)=0\), прямая \(y=c\) сама является решением. Если интеграл от \(1/g\) до \(c\) сходится, соседняя ветвь может достичь этой прямой за конечное время. Если он расходится, \(y=c\) --- горизонтальная асимптота.

\question{2}{Однородная правая часть}{Сформулируйте и поясните теорему существования и единственности
для уравнений с однородной правой частью.}
{\lectureref{10}{1604-L1630}}
\begin{result}{Теорема}
Пусть \(f\) непрерывна на интервале \(K\) и \(f(u)\ne u\) при всех \(u\in K\). Тогда для любых \(x_0\ne0\) и \(y_0/x_0\in K\) задача \(y'=f(y/x)\), \(y(x_0)=y_0\), имеет единственное локальное решение.
\end{result}
\idea{Положим \(u=y/x\), тогда \(y'=xu'+u\) и \(u'=(f(u)-u)/x\). Множители \(1/x\) и \(f(u)-u\) непрерывны, а второй отличен от нуля, поэтому применима теорема из вопроса~1.}

Если \(f(c)=c\), прямая \(y=cx\) даёт решение по обе стороны от \(x=0\). Через точки этих двух лучей единственность нужно проверять отдельно.

\clearpage
\question{3}{Уравнение высшего порядка}{Сформулируйте теорему существования и единственности для
уравнения n-го порядка, разрешенного относительно старшей
производной}
{\lectureref{4.2}{1220-L1245}}
\begin{result}{Теорема}
Рассмотрим задачу
\[
 y^{(n)}=F(x,y,y',\ldots,y^{(n-1)}),\qquad
 y^{(j)}(x_0)=a_j,\quad j=0,\ldots,n-1.
\]
Пусть в области \(G\subset\mathbb R^{n+1}\) непрерывны \(F(x,z_0,\ldots,z_{n-1})\) и все \(\partial F/\partial z_j\). Для любой начальной точки \((x_0,a_0,\ldots,a_{n-1})\in G\) найдётся \(h>0\), при котором решение существует и единственно на \([x_0-h,x_0+h]\).
\end{result}

\question{4}{Нормальная система ОДУ}{Дайте определение нормальной системы ОДУ, решения нормальной
системы ОДУ, задачи Коши для нормальной системы ОДУ,
сформулируйте теорему существования и единственности для
нормальной системы ОДУ}
{\lectureref{11}{1663-L1739}}
\textbf{Нормальная система} --- система вида \(\dot{\mathbf x}=\mathbf F(t,\mathbf x)\). Каждое из её \(n\) уравнений выражает первую производную \(\dot x^i\) через \(t,x^1,\ldots,x^n\). Пусть компоненты \(F^i\) непрерывны в области \(G\subset\mathbb R^{n+1}\).

\textbf{Решение} на промежутке \(J\) --- дифференцируемая вектор-функция \(\boldsymbol\varphi\), для которой при всех \(t\in J\) выполнены \((t,\boldsymbol\varphi(t))\in G\) и \(\dot{\boldsymbol\varphi}(t)=\mathbf F(t,\boldsymbol\varphi(t))\). В \textbf{задаче Коши} дополнительно задано \(\boldsymbol\varphi(t_0)=\mathbf x_0\).

\begin{result}{Теорема}
Если в \(G\) непрерывны также все \(\partial F^i/\partial x^j\), то через каждую точку \((t_0,\mathbf x_0)\in G\) проходит единственное решение на некотором отрезке \([t_0-h,t_0+h]\), \(h>0\).
\end{result}

\question{5}{Линейная система}{Дайте определение линейно однородной и линейно неоднородной
системы. Расскажите о применении теоремы существования и
единственности к линейной системе с непрерывной матрицей и
непрерывной вектор-функцией в правой части}
{\lectureref{11.1}{1741-L1799}}
Линейная система \(\dot{\mathbf x}=A(t)\mathbf x+\mathbf f(t)\) называется \textbf{однородной}, если \(\mathbf f\equiv0\), и \textbf{неоднородной}, если \(\mathbf f\not\equiv0\).

\begin{result}{Теорема}
Если матрица \(A\) и вектор \(\mathbf f\) непрерывны на промежутке \(I\), то при любых \(t_0\in I\), \(\mathbf x_0\in\mathbb R^n\) существует единственное решение с условием \(\mathbf x(t_0)=\mathbf x_0\). Оно определено на всём \(I\).
\end{result}
Компоненты правой части и их производные \(\partial F^i/\partial x^j=a_{ij}(t)\) непрерывны, поэтому локально применима теорема для нормальной системы. На включённом конце промежутка решение рассматривают с одной стороны. Утверждение о продолжении на весь \(I\) в лекции принимается без доказательства.

\clearpage
\question{6}{Линейная зависимость функций}{Дайте определение линейно зависимой системы функций, линейно
независимой системы вектор-функций. Сформулируйте и поясните
утверждение о связи линейной зависимости вектор-функций и их
значений. Дайте определение вронскиана n-мерных вектор-функций.}
{\lectureref{12}{1804-L1832}}
Скалярные или векторные функции \(u_1,\ldots,u_k\) \textbf{линейно зависимы} на \(M\), если \(\sum_{j=1}^k c_j u_j(t)=0\) при всех \(t\in M\) для некоторых постоянных \(c_j\), не все из которых равны нулю. Если тождество возможно лишь при \(c_1=\cdots=c_k=0\), функции \textbf{линейно независимы}.

\begin{lemma}{Лемма}
Если вектор-функции зависимы на \(M\), то их значения зависимы в каждой точке \(M\). Обратное утверждение в общем случае неверно.
\end{lemma}
\idea{В тождестве \(\sum_j c_j u_j(t)=0\) зафиксируем произвольное \(t=t_0\). Те же постоянные \(c_j\), не все нулевые, дают зависимость векторов \(u_j(t_0)\).}

\textbf{Вронскиан} \(n\) вектор-функций в \(\mathbb R^n\) --- определитель \(W(t)=\det[\,\mathbf v_1(t)\ \cdots\ \mathbf v_n(t)\,]\). Каждый столбец содержит значение одной функции в точке \(t\).

\question{7}{Вронскиан решений системы}{Сформулируйте и поясните утверждения о связи равенства нулю
определителя Вронского и линейной зависимости вектор-функций.
Сформулируйте лемму о связи равенства нулю вронскиана решений
ЛОС в некоторой точке и линейной зависимости данных функций.}
{\lectureref{12}{1834-L1879}}
Пусть \(W(t)\) --- вронскиан \(n\) вектор-функций на промежутке \(I\).
Зависимость функций влечёт \(W\equiv0\). Одной точки, где \(W(t_0)\ne0\), достаточно для независимости.

\idea{Одна и та же зависимость столбцов сохраняется при каждом \(t\). У матрицы с зависимыми столбцами определитель равен нулю.}

\begin{lemma}{Нуль вронскиана решений ЛОС}
Для \(n\) решений одной системы \(x'=A(t)x\) с непрерывной на \(I\) матрицей
из \(W(t_0)=0\) в одной точке следует, что решения зависимы на всём \(I\) и \(W\equiv0\).
Действительно, столбцы в точке \(t_0\) зависимы, поэтому их нетривиальная комбинация даёт \(z(t_0)=\sum_j c_j x_j(t_0)=0\).
По линейности \(z=\sum_j c_j x_j\) сама решает ЛОС и по единственности задачи Коши совпадает с нулевым решением на всём \(I\).
\end{lemma}
У произвольных вектор-функций \(W\equiv0\) ещё не означает зависимости.

\clearpage
\question{8}{Фундаментальная система решений}{Дайте определение фундаментальной системы решений ЛОС. Дайте
определение общего решения ЛОС. Сформулируйте и поясните
теорему об общем решении ЛОС. Дайте определение
фундаментальной матрицы ЛОС. Сформулируйте теорему о переходе
от одной фундаментальной матрицы ЛОС к другой.}
{\lectureref{13}{1883-L1962}}
\textbf{Фундаментальная система решений} (ФСР) --- это \(n\) линейно независимых
решений ЛОС порядка \(n\). \textbf{Общим решением} называют семейство решений,
которое при переборе параметров даёт все решения системы.

\begin{result}{Общее решение через ФСР}
Если \(A(t)\) непрерывна на промежутке \(I\), то пространство решений
\(x'=A(t)x\) имеет размерность \(n\). Любая ФСР \(x_1,\ldots,x_n\) --- его базис.
Поэтому каждое решение единственным образом представляется в виде
\[
 x(t)=c_1x_1(t)+\cdots+c_nx_n(t),\qquad c_j\in\mathbb R .
\]
\end{result}
\idea{У ФСР \(W(t_0)\ne0\), поэтому значения её решений образуют базис и начальный вектор единственным образом раскладывается по ним. Полученная линейная комбинация и исходное решение удовлетворяют одной задаче Коши, значит, они совпадают на всём \(I\).}

\textbf{Фундаментальная матрица} \(X(t)\) составлена из столбцов ФСР.
Все фундаментальные матрицы той же системы получаются по формуле \(Y(t)=X(t)C\),
где \(C\) --- произвольная постоянная невырожденная матрица.

\question{9}{Дифференцирование определителя}{Выпишите формулу дифференцирования определителя матрицы,
докажите её при n =2. Запишите формулу Лиувилля-Остроградского
для вронскиана решений ЛОС.}
{\lectureref{14}{1966-L1991}}
Пусть \(D(t)=\det X(t)\), а элементы \(X\) дифференцируемы.
Тогда \(D'=\sum_{i=1}^n D_i\), где в \(D_i\) только \(i\)-я строка
заменена её производной.

\textbf{Доказательство при \(n=2\).} Применим доказательство из лекции к двум строкам \((a,b)\) и \((c,d)\). Дифференцируем члены определителя и группируем их по строке, из которой взята производная.
\[
 D'=(ad-bc)'=\underbrace{a'd-b'c}_{D_1}+\underbrace{ad'-bc'}_{D_2}.
\]

\begin{result}{Формула Лиувилля--Остроградского}
Для вронскиана \(n\) решений ЛОС \(x'=A(t)x\) с непрерывной матрицей \(A\)
\[
 W(t)=W(t_0)\exp\!\left(\int_{t_0}^{t}\operatorname{tr}A(s)\,ds\right),
 \qquad \operatorname{tr}A=\sum_{i=1}^n a_{ii}.
\]
\end{result}

\clearpage
\question{10}{Вариация постоянных для системы}{Сформулируйте утверждение о связи общего решения ЛНС и ЛОС.
Расскажите о методе вариации постоянных решения ЛНС.}
{\lectureref{15}{2023-L2094}}
Пусть \(X(t)\) --- фундаментальная матрица однородной системы \(x'=A(t)x\),
а \(x_p(t)\) --- одно решение неоднородной системы \(x'=A(t)x+f(t)\) на общем промежутке непрерывности \(A\) и \(f\).
Её общее решение есть \(x=x_p+Xc\), где \(c\in\mathbb R^n\) произволен.
Действительно, разность двух решений ЛНС удовлетворяет однородной системе.

\begin{method}{Вариация постоянных}
Ищем \(x=X(t)c(t)\). Подстановка даёт \(X'c+Xc'=AXc+f\).
Так как \(X'=AX\), остаётся \(Xc'=f\), откуда \(c'=X^{-1}f\). Интегрирование даёт
\[
 x(t)=X(t)\left(c_0+\int_{t_0}^{t}X(s)^{-1}f(s)\,ds\right),
 \qquad c_0\in\mathbb R^n .
\]
\end{method}

\question{11}{Переход от уравнения к системе}{Расскажите о переходе от линейного уравнения n-го порядка к
линейной системе, сформулируйте утверждения о сохранении
линейной зависимости решений линейного уравнения и линейной
системы при данном переходе. Сформулируйте теорему об общем
решении линейного однородного уравнения (ЛОУ).}
{\lectureref{16.1}{2130-L2250}, \lectureref{16.2}{2252-L2275}}
В уравнении \(y^{(n)}+a_{n-1}y^{(n-1)}+\cdots+a_0y=f\)
с непрерывными на промежутке коэффициентами и правой частью
вводим \(x_1=y,\ x_2=y',\ldots,x_n=y^{(n-1)}\). Получаем систему
\[
 \begin{cases}
 x_k'=x_{k+1},\quad 1\le k<n,\\
 x_n'=f-a_0x_1-\cdots-a_{n-1}x_n .
 \end{cases}
\]
Первая координата решения системы даёт \(y\), остальные --- его производные.

\textbf{Линейная зависимость сохраняется в обе стороны.}
Если \(\sum_j c_jy_j\equiv0\), где постоянные \(c_j\) не все нулевые, дифференцируем это тождество до порядка \(n-1\)
и получаем те же коэффициенты зависимости для вектор-функций.
В обратную сторону достаточно взять первую координату векторного тождества.

\begin{result}{Общее решение ЛОУ}
При \(f\equiv0\) пространство решений имеет размерность \(n\).
Любые \(n\) независимых решений \(y_1,\ldots,y_n\) образуют его базис,
и общее решение имеет вид \(y=\sum_{j=1}^n c_jy_j\), \(c_j\in\mathbb R\).
\end{result}

\clearpage
\question{12}{Вронскиан скалярных функций}{Дайте определение вронскиана системы функций и расскажите о связи
линейной зависимости функций и равенства определителя Вронского
этих функций нулю. Запишите формулу Лиувилля-Остроградского для
линейных уравнений.}
{\lectureref{16.3}{2277-L2342}}
\textbf{Вронскиан функций} \(y_1,\ldots,y_n\) --- определитель
\(W(t)=\det\bigl(y_j^{(i-1)}(t)\bigr)_{i,j=1}^n\). В каждом столбце стоят
одна функция и её производные до порядка \(n-1\).

Если функции линейно зависимы на промежутке \(I\), то \(W\equiv0\).
Значит, \(W(t_0)\ne0\) хотя бы в одной точке гарантирует независимость.
Для \(n\) решений одного ЛОУ порядка \(n\) с непрерывными коэффициентами достаточно одного нуля \(W(t_0)=0\). Столбцы вронскиана решают соответствующую ЛОС, поэтому по лемме из вопроса~7 они зависимы на всём \(I\). Значит, зависимы и функции \(y_j\), а \(W\equiv0\).

\begin{result}{Формула Лиувилля--Остроградского}
Для \(n\) решений уравнения \(y^{(n)}+a_{n-1}(t)y^{(n-1)}+\cdots+a_0(t)y=0\)
\[
 W(t)=W(t_0)\exp\!\left(-\int_{t_0}^{t}a_{n-1}(s)\,ds\right).
\]
\end{result}

\question{13}{Построение уравнения и вариация постоянных}{Сформулируйте и поясните теорему о построении ЛОУ по известной
фундаментальной системе решений. Расскажите о методе вариации
постоянных решения линейного неоднородного уравнения(ЛНУ).}
{\lectureref{16.4}{2355-L2383}, \lectureref{17.2}{2473-L2523}}
\begin{result}{Построение ЛОУ по ФСР}
Пусть $y_1,\ldots,y_n\in C^n(I)$ на промежутке $I$ и их вронскиан $W(x)$ всюду отличен от нуля. Припишем $y$ к списку функций и составим вронскиан уже для $n+1$ функций. Искомое уравнение
\[
 W(y_1,\ldots,y_n,y)=0.
\]
После деления на $W(x)$ получаем ЛОУ порядка $n$ с единицей при старшей производной, непрерывными коэффициентами и заданной ФСР.
\end{result}

\idea{Разложение по последнему столбцу даёт линейное однородное уравнение с коэффициентом $W(x)\ne0$ при $y^{(n)}$. Каждая $y_j$ решает его, поскольку при $y=y_j$ два столбца совпадают, а из $W(x)\ne0$ следует независимость этих $n$ решений.}

\begin{method}{Вариация постоянных}
Пусть старший коэффициент равен единице, а справа стоит $f(x)$. Заменим постоянные при решениях ФСР функциями $c_j(x)$:
\[
 y=c_1(x)y_1+\cdots+c_n(x)y_n,\qquad
 V(x)c'(x)=(0,\ldots,0,f(x))^{\mathsf T}.
\]
В столбце $j$ матрицы $V$ стоят $y_j$ и её производные до порядка $n-1$. Так как $\det V\ne0$, система определяет $c'_j$. Интегрируем и подставляем $c_j$ в формулу для $y$.
\end{method}

\clearpage
\question{14}{Простые корни}{Дайте определение характеристического многочлена ЛОУ.
Сформулируйте и поясните теорему об общем решении ЛОУ с
постоянными коэффициентами в случае некратных корней
характеристического уравнения.}
{\lectureref{18.1}{2594-L2618}, \lectureref{19}{2732-L2776}}
\begin{definition}{Характеристический многочлен}
В ЛОУ $a_ny^{(n)}+\cdots+a_1y'+a_0y=0$ с постоянными коэффициентами и $a_n\ne0$ заменяем каждую производную $y^{(k)}$ на $\lambda^k$, а $y$ на $1$. Получаем характеристический многочлен $L(\lambda)=a_n\lambda^n+\cdots+a_1\lambda+a_0$.
\end{definition}

\begin{result}{Простые корни}
Если все $n$ корней различны, их экспоненты $e^{\lambda_1x},\ldots,e^{\lambda_nx}$ образуют ФСР. Поэтому
\[
 y=C_1e^{\lambda_1x}+\cdots+C_ne^{\lambda_nx},\qquad C_j\in\mathbb C.
\]
У вещественного уравнения пара $\alpha\pm i\beta$ даёт два вещественных решения $e^{\alpha x}\cos(\beta x)$ и $e^{\alpha x}\sin(\beta x)$.
\end{result}

\idea{Подстановка $e^{\lambda x}$ даёт $L(\lambda)e^{\lambda x}$, а для корня это ноль. При $x=0$ вронскиан экспонент равен определителю Вандермонда, который отличен от нуля для различных корней.}

\question{15}{Кратные корни и формула смещения}{Сформулируйте и поясните теорему об общем решении ЛОУ с
постоянными коэффициентами в случае кратных корней
характеристического уравнения, предварительно записав формулу
смещения}
{\lectureref{18.2}{2658-L2664}, \lectureref{20}{2782-L2879}}
\textbf{Формула смещения.}
Пусть $D=d/dx$, $L$ --- многочлен степени $n$, а $v$ имеет на $I$ производные до порядка $n$. Тогда
\[
 L(D)\bigl(e^{\lambda x}v\bigr)=e^{\lambda x}L(D+\lambda)v.
\]
\begin{result}{Кратные корни}
Корень $\lambda$ кратности $k$ даёт $k$ функций ФСР:
$e^{\lambda x},\,xe^{\lambda x},\,\ldots,\,x^{k-1}e^{\lambda x}$.
Для общего решения $L(D)y=0$ складываем по всем различным корням слагаемые $e^{\lambda x}P(x)$, где для каждого корня $P$ --- свой произвольный многочлен степени не выше $k-1$.
\end{result}

\idea{После смещения множитель $(D-\lambda)^k$ превращается в $D^k$, поэтому каждая функция $x^re^{\lambda x}$ при $r<k$ решает уравнение. В предполагаемой зависимости по очереди делим на экспоненты и берём столько производных, какова кратность очередного корня, пока не останется одна экспонента с ненулевым многочленом прежней степени, что противоречит тождеству.}

При вещественных коэффициентах для пары $\alpha\pm i\beta$ кратности $k$ берём $e^{\alpha x}\cos(\beta x)$ и $e^{\alpha x}\sin(\beta x)$ и умножаем каждую из них на $1,x,\ldots,x^{k-1}$.

\clearpage
\question{16}{Квазимногочлен и резонанс}{Дайте определение квазимногочлена. Как связаны общее решение
ЛНУ и общее решение ЛОУ? Сформулируйте теорему о частном
решении ЛНУ с постоянными коэффициентами с правой частью в виде
квазимногочлена}
{\lectureref{21}{2894-L2984}}
\begin{definition}{Квазимногочлен}
Квазимногочлен --- конечная сумма функций $e^{\mu x}P(x)$, где $\mu\in\mathbb C$, а $P$ --- многочлен.
\end{definition}

Чтобы получить общее решение ЛНУ, достаточно к одному его частному решению прибавить общее решение соответствующего ЛОУ.

\begin{result}{Выбор частного решения}
Для уравнения с постоянными коэффициентами $L(D)y=e^{\mu x}P(x)$ существует частное решение
\[
 y_{\text{ч}}=x^s e^{\mu x}Q(x),\qquad \deg Q=\deg P.
\]
Если $L(\mu)\ne0$, берём $s=0$. При резонансе $s$ равно кратности корня $\mu$. Коэффициенты $Q$ находим подстановкой. Если справа несколько слагаемых, складываем соответствующие частные решения.
\end{result}

\begin{example}{Пример квазимногочлена}
В лекции приведена сумма $e^{2x}+x^2e^{12ix}$. Для такой правой части ищут частное решение для каждого слагаемого и складывают результаты.
\end{example}

\question{17}{Интегральная формула Коши}{Сформулируйте и поясните интегральную формулу Коши нахождения
частного решения ЛНУ.}
{\lectureref{22}{2994-L3069}}
Запишем ЛНУ в виде $L[y]=f(x)$ с единицей при $y^{(n)}$. Все коэффициенты и правая часть непрерывны на промежутке $I$.

\begin{definition}{Ядро Коши}
Пусть $y_1,\ldots,y_n$ --- ФСР соответствующего ЛОУ. При фиксированном $s$ берём $K(x,s)=\sum_j c_j(s)y_j(x)$. Коэффициенты выбираем так, чтобы в точке $x=s$ производные по $x$ порядков от $0$ до $n-2$ равнялись нулю, а производная порядка $n-1$ --- единице. Для $n=1$ достаточно условия $K(s,s)=1$.
\end{definition}

\begin{result}{Интегральная формула Коши}
Для любого $\gamma\in I$ формула
\[
 Y(x)=\int_{\gamma}^{x}K(x,s)f(s)\,ds
\]
даёт частное решение с нулевыми начальными данными в точке $\gamma$.
\end{result}

\idea{При дифференцировании интеграла условия на ядро оставляют граничный член $f(x)$ только в производной порядка $n$. Остальная часть $L[Y]$ равна нулю, потому что под интегралом ядро удовлетворяет однородному уравнению.}

\end{document}
